% page 251 du fac-simile (image p-250 du scan)
% bloc 170.2 x 250.6 mm ; marge gauche 31.8 mm
\pgc{10.160mm}{0.000mm}{
\l{\cel{56}243}
\l{}
\l{\sou{ipsa}. Qua esas ya, rigoroze, to (o : ta) quon(o : quan)}
\l{\cel{3}onu jus nomis, mencionis. - L.}
\l{}
\l{\sou{-ir.}\cel{1}(gram.) Dezinenco di la infinitivo pasintala.}
\l{}
\l{\sou{irar}.(netrans.) I. (aludante vivanto). Movar su ad ula loko}
\l{\cel{3}(opoze a \textquotedbl{}venar\textquotedbl{}- \textquotedbl{}irar\textquotedbl{} signifikas for-esko di la spaco-}
\l{\cel{3}punto ube esas la persono qua parolas, ad qua onu parolas,}
\l{\cel{3}pri qua onu parolas.). II. (\sou{aludante nevivanto}).Movar su}
\l{\cel{3}(o movesar) ad ula skopo. - FIS.}
\l{}
\l{\sou{iracar}. (\sou{netrans.}) Subisar violentoza ecito, psikala,desagreabla,}
\l{\cel{3}efike da qua onu quaze eruptas (kontre ulu). - EFIS.}
\l{}
\l{\sou{\textquotedbl{}irade\textquotedbl{}}.\cel{2}Reskripto, da sultano di Konstantinopolo, pri aferi}
\l{\cel{3}min importoza kam olti qui bezonigas \textquotedbl{}hatt-i-hou-maloun\textquotedbl{} o}
\l{\cel{3}\textquotedbl{}hatti-i-oheriff\textquotedbl{}.}
\l{}
\l{\sou{iradiar}. (\sou{netrans}.) (\sou{fiziko}). Difuzar su radie. - DEFIS.}
\l{}
\l{\sou{irga.}\cel{2}Signifikas la nedetermineso maxima.}
\l{}
\l{\sou{irido}. (\sou{bot.)}\cel{2}Genero de planti monokotiledona, qua konsistas ye}
\l{\cel{3}quanto grandega de speci. - L.\cel{1}\sou{iris}.}
\l{}
\l{\sou{iridio}. (\sou{kemio)}.\cel{2}Korpo simpla, metalo ruptebla, arjento-blanka,}
\l{\cel{2}(\sou{Simbolo kemiala}\cel{1}: Ir)\cel{2}quan onu renkontras en ula platin-erci,}
\l{\cel{3}e di qua la dissolvuri esas irisea.}
\l{}
\l{\sou{irigaca}r.(\sou{trans.}) Arozar granda-quante, per kanaleti,tubi, e c.}
\l{\cel{3}- EFI.}
\l{}
\l{iriso. I. (\sou{anat.}) Parto koloroza di la okulo, membrano cirkla,}
\l{\cel{3}avan la kristalino. - II.(\sou{optiko}). Arko luma, ye la kolori}
\l{\cel{3}ek la prismo, e produktita da la refrakteso e la reflekteso}
\l{\cel{3}di la lumo-radii da la pluvo-guti. - DEFIRS.}
\l{}
\l{irisito. (\sou{patol.})\cel{2}Inflameso di la iriso.}
\l{}
\l{iritar. -(\sou{trans}.) Eventigar mikra inflameso. - DEFIS.}
\l{}
\l{ironio. I.Questiono quan Sokrates agis a la sofisti, por incitar}
\l{\cel{3}li a su-kontredici, per qui lu demonstris a li lia eroro.}
\l{\cel{3}- II. Moko-formo qua konsistas ye dicar lo inversa di to quon on}
\l{\cel{3}volas komprenigar. - DEFIRS.}
\l{}
\l{iruptar. (\sou{netrans.})\cel{2}Enirar subite e violentoze. - EFIS.}
\l{}
\l{-is.\cel{2}(\sou{gram.}) Dezinenco di la indikativo pasintala.}
\l{}
\l{iskiono. (\sou{anat.)}\cel{2}Parto infra di la osti hanchala,ube inkastresas}
\l{\cel{3}la osto di la kruro.}
\l{}
\l{\sou{islamo}. I. Religio e civilizeso-formo di la Mohamedisti. - II. En-}
\l{\cel{2}semblo ek la landi quin karakterizas ta religio e ta civilizeso-}
\l{\cel{2}formo.}
}
